\documentclass[12pt,a4paper]{article}

\usepackage[utf8]{inputenc}
\usepackage[T1]{fontenc}
\usepackage[margin=2.5cm]{geometry}
\usepackage{parskip}

\title{A Topological Loss Function for Deep-Learning Based Image Segmentation Using Persistent Homology, IEEE TPAMI}
\author{}
\date{}

\begin{document}
\maketitle

\section*{1. Place of origin}
    The article is printed in the international journal
    IEEE Transactions on Pattern Analysis and Machine Intelligence
    (IEEE TPAMI), published by the IEEE Computer Society. Has Q1 quartile.

\section*{2. Time of origin}
    The article is dated december 2022 (volume 44, number 12).

\section*{3. Author}
    The article is written by a group of authors - James Clough,
    Nicholas Byrne, Ilkay Oksuz, Veronika Zimmer, Julia Schnabel,
    and Andrew King, mostly affiliated in London and one from Istanbul and one from German.

\section*{4. Theme / Topic}
    The article deals with training convolutional neural networks (CNN)
    for image and volume segmentation. It addresses the problem that
    standard segmentation loss functions ignore the large-scale shape,
    or topology, of the segmented object, and proposes a way to
    incorporate prior topological knowledge into training using
    persistent homology, a concept from topological data analysis.

\section*{5. Main idea / Aim of the article}
    The main idea is that pixel-wise loss functions, such as
    cross-entropy or the Dice loss, only measure overlap between the
    predicted and ground-truth segmentations and cannot capture whether
    its large-scale structure is topologically correct. The aim of the
    authors is to provide a differentiable topological loss function
    that specifies the desired Betti numbers of an object - the number
    of connected components, loops, and cavities - and to show that
    driving the network towards the correct topology also improves
    pixel-wise accuracy, even without ground-truth labels for every
    case.

\section*{6. Contents of the article}
    The article can be divided into four parts.
    The first part deals with the introduction and related work,
    explaining why pixel-wise losses fail to capture topology.
    The second covers the theory: persistent homology of cubical
    complexes, Betti numbers, barcode diagrams, and the topological loss
    function $L_{topo}$, applied either as post-processing or in a
    semi-supervised framework.
    The third touches upon four experiments: denoising MNIST digits;
    segmenting the hearts myocardium from MRI data; the public ACDC
    challenge dataset; and segmenting the placenta from 3-D ultrasound.
    The fourth part includes the discussion and conclusion, addressing
    the methods limitations and further applications.

    The article is written mostly in the first person plural (like :we)
    and includes equations, result tables, and figures such as barcode
    diagrams. The authors start by stating that pixel-wise losses cannot
    capture topological correctness, then describe the theory and derive
    the loss function. According to the text, across all four
    experiments the topological loss consistently improved both the
    Dice score and the proportion of topologically correct
    segmentations, especially on harder tasks. The authors conclude
    that persistent homology is a viable tool for embedding topological
    prior knowledge into segmentation training.

\section*{7. Personal opinion / impression of the article}
    I found the article interesting and useful, since it presents a
    genuinely new approach to combining topology with machine learning.
    Because topology is, by its nature, concerned with transformations
    of shapes and connectivity patterns, it seems logical that it finds
    application in computer vision, and in segmentation in particular. I
    believe this direction - applying persistent homology to
    segmentation - is very likely to become the topic of my own
    masters thesis, aiming at either faster inference, faster training,
    or better accuracy than classical pixel-wise loss functions.

\section*{8. Personal view on the topic / idea / problem}
    I share the authors view that pixel-wise loss functions alone are
    insufficient whenever the global shape of an object matters. At the
    same time, I see the problem in a slightly different way: computing
    persistent homology adds computational cost, on the
    order of seconds per batch or per volume, which makes the method
    unsuitable for real-time segmentation. To my mind, this is the main
    bottleneck of the approach, and improving the speed of the
    computation is the most important open direction for future work,
    including a possible direction for my own research. I dont quite
    agree that computational optimisation should be left outside the
    scope of the paper: without it, topological loss functions are
    unlikely to be widely adopted in real clinical or industrial
    pipelines.

\end{document}
